\begin{table}[H]
\centering
\footnotesize
\caption{Distance from the source data and the gain from recalibration in the nine development regions (exploratory, computed after the confirmation run). Outside: share of the region's FUA-months outside the area of applicability of the other eight development regions; DI: median dissimilarity index. Gain: mean paired RMSE reduction (\%) of source-affine relative to direct transfer over the region's 24 pseudo-transfers (five labelled FUAs, three to twelve months); negative values favour direct transfer. Wins: pseudo-transfers in which source-affine is more accurate.}
\label{tab:adapt-distance}
\setlength{\tabcolsep}{4pt}
\begin{tabular}{@{}lrrrrr@{}}
\toprule
Region & FUAs & Outside (\%) & DI & Gain (\%) & Wins \\
\midrule
North China Plain (Beijing, Jinan, Qingdao) & 40 & 1.6 & 0.17 & $-$0.4 & 12/24 \\
Shanxi, Henan and Shaanxi (Taiyuan, Xi'an) & 45 & 2.7 & 0.19 & $-$19.1 & 1/24 \\
Central China and Fujian (Wuhan, Changsha, Xiamen) & 44 & 4.8 & 0.17 & $-$14.9 & 2/24 \\
Sichuan Basin and Guizhou (Chengdu, Chongqing) & 38 & 8.7 & 0.21 & $-$21.3 & 1/24 \\
South China (Guangzhou, Nanning, Haikou) & 42 & 10.9 & 0.20 & $-$4.3 & 11/24 \\
Liaoning and north Hebei (Shenyang, Tangshan) & 27 & 11.2 & 0.22 & $-$13.6 & 1/24 \\
Yunnan and Tibet (Kunming, Lhasa) & 17 & 14.8 & 0.22 & 26.4 & 23/24 \\
North-east (Harbin, Changchun) & 27 & 17.4 & 0.23 & $-$42.3 & 0/24 \\
Xinjiang (Urumqi, Karamay) & 14 & 100.0 & 0.99 & 18.5 & 23/24 \\
\bottomrule
\end{tabular}
\end{table}
